% Generated draft: requires mathematical and source review.
\documentclass[11pt]{article}
\usepackage[margin=25mm]{geometry}
\usepackage{fontspec}
\setmainfont{Latin Modern Roman}
\usepackage{amsmath,amssymb,mathtools,unicode-math}
\setmathfont{Latin Modern Math}
\usepackage{enumitem,needspace,xurl,hyperref}
\setlength{\parindent}{0pt}
\setlength{\parskip}{5pt}
\setlength{\emergencystretch}{3em}
\newcommand{\N}{\mathbb{N}}
\newcommand{\Z}{\mathbb{Z}}
\newcommand{\Q}{\mathbb{Q}}
\newcommand{\R}{\mathbb{R}}
\newcommand{\C}{\mathbb{C}}
\newcommand{\F}{\mathbb{F}}
\newcommand{\A}{\mathbb{A}}
\newcommand{\PP}{\mathbb{P}}
\newcommand{\E}{\mathbb{E}}
\newcommand{\eps}{\varepsilon}
\newcommand{\dd}{\,\mathrm{d}}
\begin{document}
\section*{UnsolvedMath ID: 700757}
\textbf{Source problem code:} catalogue entry
\subsection{Definitions and mathematical statement}
Let \(p\) be a prime. A pro-\(p\) group is a topological group isomorphic to an inverse limit of finite \(p\)-groups. Such a group \(G\) is topologically finitely generated if there is a finite subset \(S\subseteq G\) such that \(\overline{\langle S\rangle}=G\). It is torsion-free if, for every \(g\in G\) and every integer \(n\geq 1\), the equality \(g^n=1\) implies \(g=1\). It is soluble if its derived series, defined by \(G^{(0)}=G\) and \(G^{(i+1)}=[G^{(i)},G^{(i)}]\), reaches the trivial subgroup after finitely many steps.

In the usual first-order language of groups \(\mathcal L_{\mathrm{grp}}=\{\cdot,{}^{-1},1\}\), let \(\operatorname{Th}_{\mathcal L_{\mathrm{grp}}}(G)\) be the set of all sentences true in the underlying abstract group \(G\). This theory is decidable if an algorithm determines, for every sentence in this language, whether it belongs to that set. An analytic pro-\(p\) group means a pro-\(p\) group admitting a finite-dimensional \(\Q_p\)-analytic manifold structure inducing its given topology, with analytic multiplication and inversion.

\textbf{Conjecture (A. G. Myasnikov and V. N. Remeslennikov).} For every prime \(p\), every topologically finitely generated, soluble, torsion-free pro-\(p\) group \(G\) whose elementary theory is decidable is an analytic pro-\(p\) group.

\textit{Formulation qualification.} The supplied statement does not specify the language of the elementary theory or explicitly distinguish topological from abstract finite generation. The formulation above uses the usual pure group language and topological finite generation; these interpretations require confirmation against the original source.
\subsection{Short English statement}
The conjecture asks whether every finitely generated soluble torsion-free pro-\(p\) group with decidable elementary theory must be a finite-dimensional \(p\)-adic analytic group. Here finite generation is interpreted topologically and elementary theory in the pure group language, subject to source confirmation.
\subsection{Sources}
\begin{itemize}
\item Local UnsolvedMath registry, supplied record with permanent ID \(700757\), exact source code \texttt{catalogue entry}, titled \emph{Kourovka Notebook Problem 11.65}. The record supplies the status \texttt{open}; this status is not independently verified.
\item \emph{The Kourovka Notebook (Unsolved Problems in Group Theory)}, problem 11.65, attributed to A. G. Myasnikov and V. N. Remeslennikov, proposed in the 11th issue (1990). The supplied background also identifies the 21st issue edition (2026), arXiv:1401.0300v48. The collection description names E. I. Khukhro and V. D. Mazurov as editors. These bibliographic details are supplied attribution, not independently verified source consultation.
\end{itemize}
\end{document}
